% ==============================================================================
% abstract.tex
% ==============================================================================

\chapter*{Abstract}
\addcontentsline{toc}{chapter}{Abstract}

\todo[inline]{Write abstract once the work is complete (recommended: 250--300 words).}

Urban air quality management increasingly demands spatially resolved emission
inventories that capture the heterogeneity of pollutant sources at the city
scale. The \urbem{} tool provides a hybrid, automated framework to downscale
pan-European \cams emission inventories ($\sim \qtyproduct[product-units=single]{6 x 6}{\km\squared}$)
to a high-resolution grid of \qtyproduct[product-units=single]{1 x 1}{\km\squared}, using publicly
available geospatial proxies (OpenStreetMap road network, CORINE Land Cover,
Global Human Settlement population data, E-PRTR point sources).

This master's thesis presents the upgrade of \urbem{} to ingest the latest
\cams{} datasets produced in the RI-URBANS project (v6.1, base year 2019),
incorporating improved road transport emissions (bottom-up COPERT methodology,
updated spatial allocation via OpenTransportMap) and novel ultrafine particle
(UFP) number inventories. The upgraded pipeline is applied over two \textit{Horizon
Europe} Metronome Mission Cities in Greece — Ioannina and Kozani — to produce
city-scale emission inputs for the \episode{} chemistry transport model.

Model predictions are evaluated against in-situ observations, and the effects
of the updated inventory on modelled concentration fields are quantified. The
framework is further used to assess the co-benefits on air quality of urban
climate mitigation scenarios targeting the building and road transport sectors.

\bigskip
\noindent\textbf{Keywords:} urban air quality, emission inventory, downscaling,
CAMS-REG, UrbEm, ultrafine particles, EPISODE-CityChem, Ioannina, Kozani,
Metronome project.
